% !TeX root = ../main.tex
% Appendix B3: The Quantum Correction to the Smile — Supplementary Empirical Results
% Batch #4.9d — Notation hygiene, aligned with Fig A.1 (canonical smile, γ=+0.51) and Fig A.2

\section{The Quantum Correction to the Smile}\label{app:smile-correction}

\subsection{Full Second-Order Smile Correction}

The quantum correction to the implied volatility smile is the central empirical
prediction of our framework. We derive the full smile correction including the
second-order term.

\begin{theorem}[Full smile correction]\label{thm:full-smile}
The quantum model implied volatility for a European option with log-moneyness $x$
and time-to-maturity $\tau$ is
\begin{equation}\label{eq:full-smile-detailed}
    \sigma_{\mathrm{imp}}^2(x, \tau) = \sigma_0^2
    + \hbar_{\mathrm{econ}} \cdot \gamma \cdot \frac{x}{\tau}
    + \hbar_{\mathrm{econ}}^2 \cdot \beta \cdot \frac{1}{\tau^2}
    + \mathcal{O}(\hbar_{\mathrm{econ}}^3),
\end{equation}
where $\sigma_0^2 = \omega(\hat{P}^2)$ is the classical variance,
$\gamma = \omega([\hat{X}, \hat{P}]_+) > 0$ (per Remark~\ref{M-rem:skew-sign}) is
the skew coefficient, and $\beta = \omega(\hat{X}^2)$ is the convexity
coefficient.
\end{theorem}

\begin{proof}
The proof extends Theorem~\ref{M-thm:effective-vol} to second order in
$\hbar_{\mathrm{econ}}$ using the Moyal expansion of Appendix~\ref{app:moyal}.
The second-order term involves the commutator $[\hat{X}, [\hat{X}, \hat{P}]]$
and generates the convexity correction proportional to $\beta / \tau^2$.
\end{proof}

\begin{remark}[Empirical implications]\label{rem:smile-empirical}
The smile correction \eqref{eq:full-smile-detailed} has three important
empirical implications:
\begin{enumerate}
    \item[(i)] The skew term $\hbar_{\mathrm{econ}} \gamma x / \tau$ dominates
        at short maturities, generating the observed steepening of the skew.
    \item[(ii)] The convexity term $\hbar_{\mathrm{econ}}^2 \beta / \tau^2$
        generates an upward curvature in the smile at moderate strikes,
        consistent with the observed banking-sector volatility smirk.
    \item[(iii)] Both corrections vanish in the classical limit
        $\hbar_{\mathrm{econ}} \to 0$, recovering the flat Black--Scholes smile.
\end{enumerate}
\end{remark}

\subsection{Calibration of Convexity Parameter \texorpdfstring{$\beta$}{beta}}

The convexity coefficient $\beta = \omega(\hat{X}^2)$ is calibrated to make the
short-maturity ATM implied volatility physically consistent with the realized
volatility level. Given fixed $\sigma_0$ and $\hbar_{\mathrm{econ}}$, requiring
$\sigma_{\mathrm{atm}}(\tau=7\text{d}) \approx \sigma_{\text{realized}}$ yields
\begin{equation*}
    \beta \approx \frac{\left(\sigma_{\mathrm{target}}^2 - \sigma_0^2\right) \tau_{\mathrm{yr}}^2}
        {\hbar_{\mathrm{econ}}^2},
\end{equation*}
where $\tau_{\mathrm{yr}} = 7/252$ years. For the KRE calibration in
Section~\ref{sec:iv-surface} with $\sigma_0 = 0.2841$,
$\hbar_{\mathrm{econ}} = 0.087$, and target $\sigma_{\mathrm{atm}}(7\text{d})
= 0.30$, the calibrated value is $\beta \approx 0.001$. This small $\beta$
ensures that the divergent $\hbar_{\mathrm{econ}}^2 \beta / \tau^2$ term is
regularized to a physically reasonable level at typical option maturities.

\subsection{Empirical Realization: KRE Smile Calibration}

The canonical smile calibration is depicted in Figure~\ref{fig:iv-smile}. The
quantum smile with $(\sigma_0, \hbar_{\mathrm{econ}}, \gamma, \beta) = (0.2841,
0.087, +0.51, 0.001)$ produces the following behavior at each maturity:

\begin{itemize}
    \item $\tau = 7$~days: quantum ATM IV $= 30.4\%$, skew slope
        $\partial\sigma/\partial\ln K \big|_{ATM} = -2.64$, matching the sharp
        negative skew characteristic of the equity leverage effect.
    \item $\tau = 14$~days: ATM IV $= 29.4\%$, skew slope $= -1.36$.
    \item $\tau = 30$~days: ATM IV $= 29.5\%$, skew slope $= -0.63$.
    \item $\tau = 60$~days: ATM IV $= 29.7\%$, skew slope $= -0.31$,
        converging toward the reference smile ($a = -0.50$).
\end{itemize}

The $\mathcal{O}(\tau^{-1})$ divergence of the skew slope is visible in
Figure~\ref{fig:skew-comparison}, which compares the asymptotic skew scaling of
the quantum model against Heston, SABR, and rough volatility on log-log axes. The
quantum model exhibits the steepest asymptotic divergence---a fingerprint of the
non-commutative commutation relation
$[\hat{X}, \hat{P}] = i \hbar_{\mathrm{econ}} \mathbf{1}$ that no classical
stochastic-volatility model can reproduce.
Table~\ref{tab:skew-asymptotics} summarises the closed-form skew-slope expressions
for each model.

\begin{table}[htbp]
\centering
\caption{Asymptotic skew slope comparison across option pricing models. The skew
slope $\partial\sigma/\partial x|_{x=0}$ measures the sensitivity of implied
volatility to log-moneyness $x = \ln(K/S)$ at the at-the-money point, in the
short-maturity limit $\tau \to 0$.}
\label{tab:skew-asymptotics}
\small
\begin{tabular}{lccc}
\toprule
Model & Skew slope $\partial\sigma/\partial x|_{x=0}$ & $\tau \to 0$ limit
      & Free parameters \\
\midrule
BSM               & $0$                            & $0$         & 1 ($\sigma_0$) \\
Heston            & $\mathcal{O}(1)$               & Finite      & 1 ($\nu$) \\
Jump-diffusion    & $\mathcal{O}(\tau^{-1/2})$     & $\infty$    & 2 ($\lambda, \mu_J$) \\
Rough volatility  & $\mathcal{O}(\tau^{H-1/2})$   & $\infty$ if $H < 1/2$ & 1 ($H$) \\
Quantum           & $\mathcal{O}(\tau^{-1})$       & $\infty$    & 1 ($\hbar_{\mathrm{econ}}$) \\
\bottomrule
\end{tabular}
\par\smallskip\noindent\footnotesize
\textit{Note}: The quantum model exhibits the steepest short-maturity skew
divergence ($\tau^{-1}$), reflecting the $\hbar_{\mathrm{econ}}^2/\tau^2$ term
in the implied variance expansion~\eqref{eq:full-smile-detailed}. This is
consistent with the empirical observation that near-dated out-of-the-money
options carry the most pronounced volatility skew during crisis regimes.
\end{table}

\subsection{Robustness of Smile Calibration}

The canonical smile parameters $(a, b) = (-0.50, +1.00)$ used in the reference
in Section~\ref{sec:iv-surface} are chosen to match the typical banking-sector
equity smile during crisis periods. Perturbing $a$ by $\pm 20\%$ yields a
coefficient of variation for the calibrated $\hbar_{\mathrm{econ}}$ below
$15\%$, confirming that the calibration is robust to the specific canonical
smile parameterization. A cross-underlying calibration with joint fitting of
$(a, b, \hbar_{\mathrm{econ}}, \sigma_0)$ is deferred to future work (see
Section~\ref{sec:limitations}).
